\documentclass[14pt,a4paper]{article}

% --- Подключение общего стиля из родительской папки ---
\makeatletter
\def\input@path{{../}}
\makeatother
\usepackage{common}

% --- Папка для картинок этого документа ---
\graphicspath{{fig/}}

\begin{document}

% ===================== ТИТУЛЬНЫЙ ЛИСТ =====================
\begin{titlepage}
\begin{center}
МИНИСТЕРСТВО НАУКИ И ВЫСШЕГО ОБРАЗОВАНИЯ РОССИЙСКОЙ ФЕДЕРАЦИИ\\
ФЕДЕРАЛЬНОЕ ГОСУДАРСТВЕННОЕ АВТОНОМНОЕ ОБРАЗОВАТЕЛЬНОЕ УЧРЕЖДЕНИЕ\\
ВЫСШЕГО ОБРАЗОВАНИЯ\\
«САМАРСКИЙ НАЦИОНАЛЬНЫЙ ИССЛЕДОВАТЕЛЬСКИЙ УНИВЕРСИТЕТ\\
ИМЕНИ АКАДЕМИКА С. П. КОРОЛЕВА»\\
(Самарский университет)\\[12pt]
Кафедра \underline{\hspace{6cm}}
\end{center}

\vfill

\begin{center}
РУКОВОДСТВО ПОЛЬЗОВАТЕЛЯ\\
языка программирования \texttt{BasicLite}\\
(вариант 3-2-1-A)
\end{center}

\vfill

\begin{flushright}
\begin{tabular}{l}
Студент группы \underline{\hspace{2.5cm}} \\
\underline{\hspace{5cm}} \\
Руководитель \underline{\hspace{2.5cm}} \\
\underline{\hspace{5cm}}
\end{tabular}
\end{flushright}

\vfill

\begin{center}
Самара 2024
\end{center}
\end{titlepage}

\structtitle{Содержание}
\tableofcontents

% ===================== ВВЕДЕНИЕ =====================
\structtitle{Введение}

\texttt{BasicLite} — упрощённый императивный язык с BASIC-подобным
синтаксисом, строгой типизацией и циклом с постусловием. Язык
предназначен для учебных целей и поддерживает редактирование,
пошаговое выполнение и отладку.

% ===================== 1 БАЗОВЫЙ СИНТАКСИС =====================
\section{Базовый синтаксис}

\subsection{Регистр}

Ключевые слова не чувствительны к регистру: \texttt{PRINT},
\texttt{print}, \texttt{Print} эквивалентны. Идентификаторы также
не чувствительны к регистру (\texttt{X} и \texttt{x} — одна
переменная).

\subsection{Комментарии}

Однострочный комментарий начинается с \texttt{'}:
\begin{lstlisting}[language={},numbers=none]
' это комментарий
PRINT "hello"   ' и это тоже
\end{lstlisting}

\subsection{Разделители операторов}

Операторы разделяются концом строки. Точка с запятой не используется.
Пустые строки игнорируются.

% ===================== 2 ТИПЫ ДАННЫХ =====================
\section{Типы данных}

Язык поддерживает три типа, перечисленные в таблице~\ref{tab:types}.

\begin{table}[H]
\caption{Типы данных языка}
\label{tab:types}
\centering
\begin{tabular}{@{}lll@{}}
\toprule
Тип & Пример литерала & Описание \\
\midrule
INTEGER & \texttt{42}      & целое число \\
REAL    & \texttt{3.14}    & вещественное число \\
STRING  & \texttt{"hello"} & строка \\
\bottomrule
\end{tabular}
\end{table}

Правила строгой типизации:
\begin{itemize}
  \item преобразование INTEGER $\to$ REAL разрешено;
  \item преобразование REAL $\to$ INTEGER запрещено;
  \item STRING не участвует в арифметике;
  \item смешивать STRING и числа в одной операции нельзя.
\end{itemize}

% ===================== 3 ПЕРЕМЕННЫЕ =====================
\section{Переменные}

\subsection{Объявление}

\begin{lstlisting}[language={},numbers=none]
DIM x AS INTEGER
DIM y AS REAL
DIM s AS STRING
DIM a, b AS INTEGER
\end{lstlisting}

\subsection{Присваивание}

Оператор присваивания — \texttt{:=}:
\begin{lstlisting}[language={},numbers=none]
x := 10
y := 3.5
s := "hello"
x := x + 1
\end{lstlisting}

% ===================== 4 ВЫРАЖЕНИЯ =====================
\section{Выражения и операции}

\subsection{Арифметические операции}

\texttt{+}, \texttt{-}, \texttt{*}, \texttt{/}. Приоритет:
\texttt{*}, \texttt{/} выше, чем \texttt{+}, \texttt{-}. Скобки
\texttt{()} меняют порядок.

\subsection{Операции сравнения}

\texttt{>}, \texttt{<}, \texttt{>=}, \texttt{<=}, \texttt{=},
\texttt{<>}. Результат — INTEGER (0 или 1).

\subsection{Логические операции}

\texttt{AND}, \texttt{OR}, \texttt{NOT}. Приоритет:
\texttt{NOT} > \texttt{AND} > \texttt{OR}.

\subsection{Пример}

\begin{lstlisting}[language={},numbers=none]
DIM a, b AS INTEGER
a := (5 + 3) * 2
b := 10 / 3          ' ошибка: / даёт REAL, нельзя в INTEGER
\end{lstlisting}

% ===================== 5 ВВОД-ВЫВОД =====================
\section{Операторы ввода-вывода}

\subsection{PRINT}

\begin{lstlisting}[language={},numbers=none]
PRINT "hello"
PRINT x
PRINT "x =", x, "y =", y
PRINT
\end{lstlisting}

\subsection{INPUT}

\begin{lstlisting}[language={},numbers=none]
INPUT x
INPUT a, b
\end{lstlisting}

% ===================== 6 УПРАВЛЯЮЩИЕ КОНСТРУКЦИИ =====================
\section{Управляющие конструкции}

\subsection{Условный оператор}

\begin{lstlisting}[language={},numbers=none]
IF x > 0 THEN
    PRINT "positive"
ELSE
    PRINT "non-positive"
ENDIF
\end{lstlisting}

\subsection{Цикл с постусловием}

Тело выполняется хотя бы один раз; проверка — после тела.
\begin{lstlisting}[language={},numbers=none]
DO
    PRINT i
    i := i + 1
LOOP WHILE i <= 10

DO
    PRINT i
    i := i + 1
LOOP UNTIL i > 10
\end{lstlisting}

\subsection{Цикл со счётчиком}

\begin{lstlisting}[language={},numbers=none]
FOR i = 1 TO 10
    PRINT i
NEXT

FOR i = 10 TO 1 STEP -1
    PRINT i
NEXT
\end{lstlisting}

\subsection{Множественный выбор}

\begin{lstlisting}[language={},numbers=none]
SELECT CASE x
CASE 1, 2, 3:
    PRINT "small"
CASE 4 TO 10:
    PRINT "medium"
CASE ELSE:
    PRINT "large"
ENDSELECT
\end{lstlisting}

\subsection{Безусловный переход}

\begin{lstlisting}[language={},numbers=none]
GOTO done
PRINT "skipped"
done:
PRINT "end"
\end{lstlisting}

% ===================== 7 ФУНКЦИИ =====================
\section{Функции}

\subsection{Определение}

\begin{lstlisting}[language={},numbers=none]
FUNCTION square(x AS INTEGER) AS INTEGER
    RETURN x * x
ENDFUNCTION
\end{lstlisting}

\subsection{Вызов}

\begin{lstlisting}[language={},numbers=none]
DIM y AS INTEGER
y := square(5)
PRINT y
CALL square(3)
\end{lstlisting}

% ===================== 8 ПРИМЕР =====================
\section{Пример программы}

\begin{lstlisting}[language={},numbers=none]
' Сумма чисел от 1 до N
DIM N, I, S AS INTEGER

PRINT "Enter N:"
INPUT N

I := 1
S := 0
DO
    S := S + I
    I := I + 1
LOOP WHILE I <= N

PRINT "Sum =", S

IF S > 100 THEN
    PRINT "Large sum"
ELSE
    PRINT "Small sum"
ENDIF
\end{lstlisting}

% ===================== 9 РАБОТА С ИНТЕРПРЕТАТОРОМ =====================
\section{Работа с интерпретатором}

\subsection{Запуск}

Порядок запуска:
\begin{enumerate}
  \item запустить \texttt{python main.py};
  \item нажать «Загрузить» и выбрать файл \texttt{.bas} либо набрать
        программу в редакторе;
  \item нажать «Запустить» для полного выполнения.
\end{enumerate}

% ФОТО: главное окно после запуска. Положите в manual/fig/run.png
\begin{figure}[H]
  \centering
  \fbox{\parbox[c][7cm][c]{0.9\textwidth}{\centering
    \textbf{[МЕСТО ДЛЯ ФОТО]}\\[4pt]
    Главное окно после запуска программы}}
  \caption{Запуск программы}
  \label{fig:run}
\end{figure}

\subsection{Пошаговое выполнение}

Кнопка «Шаг» выполняет одну инструкцию ПОЛИЗа и обновляет панель
ПОЛИЗа (текущая инструкция подсвечена), таблицу переменных и стек
(в строке статуса).

% ФОТО: пошаговое выполнение. Положите в manual/fig/step.png
\begin{figure}[H]
  \centering
  \fbox{\parbox[c][7cm][c]{0.9\textwidth}{\centering
    \textbf{[МЕСТО ДЛЯ ФОТО]}\\[4pt]
    Пошаговое выполнение: подсветка ПОЛИЗа, стек, переменные}}
  \caption{Пошаговое выполнение}
  \label{fig:step}
\end{figure}

\subsection{Точки останова}

Порядок работы с точками останова:
\begin{enumerate}
  \item установить курсор на нужную инструкцию ПОЛИЗа;
  \item нажать «Точка останова» — инструкция помечается;
  \item нажать «Запустить» — выполнение остановится на этой
        инструкции;
  \item для продолжения — снова «Запустить» или «Шаг».
\end{enumerate}

% ФОТО: точка останова. Положите в manual/fig/bp.png
\begin{figure}[H]
  \centering
  \fbox{\parbox[c][7cm][c]{0.9